\chapter{Einleitung}\label{ch:einleitung}

\section{Motivation}

% TODO V14+: Motivation aus Habich-Sprechstunden + 33-Paper-Korpus
% Aufhaenger: Cache-Awareness in Trie-basierten Index-Strukturen ist
% historisch in der Forschung weit gestreut, aber nicht systematisch
% permutiert + ausgemessen worden.

\section{Forschungsfragen}

\begin{itemize}
    \item \textbf{F1}: Welche Algorithmus-Bausteine eines Cache-bewussten
          Trie-Index lassen sich in einer 11-Achsen-Permutationsmatrix
          systematisch kombinieren?
    \item \textbf{F2}: Welche Permutationen erzielen die besten
          Latenz/Throughput-Werte unter YCSB-A bis YCSB-F?
    \item \textbf{F3}: Wie verhaelt sich PRT-ART (Pruefling) im Vergleich
          zum Stand der Technik (8 Tier-1-SOTA + 22 Tier-2/3-SOTA)?
\end{itemize}

\section{Beitraege}

\begin{itemize}
    \item Drei-Schichten-Architektur (CacheEngine $\to$ ExecutionEngine $\to$
          SearchEngine) mit ABI-stabilen C++23-Modul-Interfaces.
    \item 30 SOTA-Algorithmus-Profile als persistierte XML-Konfigurationen.
    \item PRT-ART als Pruefling mit 8 eigenen Bausteine-Schichten.
    \item Mess-Driver mit Defined/Full-Mode + Profile-aware Workload-Routing.
\end{itemize}

\section{Aufbau der Arbeit}

Kapitel~\ref{ch:sota} fasst den 33-Paper-Stand der Technik zusammen.
Kapitel~\ref{ch:architektur} stellt die Drei-Schichten-Architektur vor.
Kapitel~\ref{ch:implementation} dokumentiert die Implementation in den drei
Repositorien (cache-engine, prt-art, Diplomarbeit-Code). Kapitel~\ref{ch:messverfahren}
beschreibt das Mess-Verfahren. Kapitel~\ref{ch:auswertung} praesentiert die
Mess-Resultate. Kapitel~\ref{ch:fazit} schliesst mit Fazit und Ausblick.
